7. Sea \( V \) un \( \mathbb{R} \)-espacio vectorial con un producto escalar \( \langle-,-\rangle \) positivo definido y \( W \) el subespacio generado por el conjunto ortonormal \( \left\{v_{1}, v_{2}, \ldots, v_{k}\right\} \).
(a) Demuestra la Identidad de Parseval: Dados \( u \in W \) y \( w \in V \), entonces
\[
\langle u, w\rangle=\sum_{i=1}^{k}\left\langle u, v_{i}\right\rangle\left\langle v_{i}, w\right\rangle
\]
(b) Demuestra la Desigualdad de Bessel: Para cualesquiera \( v \in V \)
\[
\sum_{i=1}^{k}\left\langle u, v_{i}\right\rangle^{2} \leq\|u\|^{2}
\]

La igualdad se satisface si \( u \in W \). \\

\begin{solution}

\textbf{Parte (a): Demostración de la Identidad de Parseval.} \\
Dado que \( u \in W \), \( u \) se puede expresar como una combinación lineal de los vectores ortonormales \( v_1, v_2, \ldots, v_k \). Así, podemos escribir:
\[
   u = \sum_{i=1}^{k} c_i v_i
   \]
donde \( c_i = \langle u, v_i \rangle \) debido a la propiedad de ortonormalidad.\\
Al calcular el producto escalar \(\langle u, w \rangle\), tenemos: 
\[
   \langle u, w \rangle = \left\langle \sum_{i=1}^{k} c_i v_i, w \right\rangle = \sum_{i=1}^{k} c_i \langle v_i, w \rangle
   \]
Sustituyendo \( c_i = \langle u, v_i \rangle \), obtenemos:
\[
   \langle u, w \rangle = \sum_{i=1}^{k} \langle u, v_i \rangle \langle v_i, w \rangle
   \]
\textbf{Parte (b): Demostración de la Desigualdad de Bessel}.\\
Consideremos la proyección ortogonal de \( v \) sobre \( W \), denotada por \( v_W \). La proyección \( v_W \) se define como:
\[
u_W = \sum_{i=1}^{k} \langle u, v_i \rangle v_i
\]
De aquí se sigue que.
\[
\|u_W\|^2 = \left\langle \sum_{i=1}^{k} \langle u, v_i \rangle v_i, \sum_{j=1}^{k} \langle u, v_j \rangle v_j \right\rangle = \sum_{i=1}^{k} \langle u, v_i \rangle^2
\]
Ahora, descomponemos \( u \) como la suma de su proyección \( u_W \) y una componente ortogonal \( u_\perp \) a \( W \).
\[
u = u_W + u_\perp
\]
Dado que \( u_W \) es ortogonal a \( u_\perp \), por el teorema de Pitágoras, tenemos:
\[
\|u\|^2 = \|u_W\|^2 + \|u_\perp\|^2 \geq \|u_W\|^2 = \sum_{i=1}^{k} \langle u, v_i \rangle^2
\]
Esto demuestra que:
\[
\sum_{i=1}^{k} \langle u, v_i \rangle^2 \leq \|u\|^2
\]
La igualdad se cumple si y solo si \( u_\perp = 0 \), es decir, si \( u \in W \), en cuyo caso \( u = u_W \).


\end{solution}